\section{Open questions (5)}

The following are questions this replication does \emph{not} answer but
which flow directly from the work; each includes an evidentiary basis
and a concrete next-step probe. Source of truth:
\texttt{report/open\_questions.json}.

\subsection*{Q1. Do the paper's verdicts survive on the exact reduced Hamiltonians of the cited experiments?}
\textbf{Basis.} Our $|S|$ differs from Table~I by $3{\times}$--$9{\times}$
because we use default OpenFermion transforms rather than Hempel/Peruzzo/Nam's
exact experimental reductions. The verdict is invariant under the reductions
that trim $S$, but a subtle representation-dependent flip is not ruled out
at the boundary (e.g., $\mathrm{HeH}^+$ with $CD_0{=}0.38$).\\
\textbf{Next steps.} Pull the exact Pauli lists from each experiment's SI
or code; rerun Theorem~3; confirm $|S|$ and verdict match; document any
flip and which reduction step introduces it. $\sim$1--2 CPU-hours per
molecule; no paid APIs.

\subsection*{Q2. Can one actually construct and evaluate the noncontextual hidden-variable model for a non-contextual VQE instance, and does it reproduce the VQE trajectory?}
\textbf{Basis.} The paper's punchline is that non-contextual $\Rightarrow$
classically simulable. We verified only the classifier; we did not build a
classical simulator. Kirby-Love (2020) provides the construction.\\
\textbf{Next steps.} Implement the Kirby-Love 2020 classical simulator for
2q $\mathrm{H}_2$, 4q $\mathrm{H}_2$ JW/BK; compare energies to VQE at 5--10
bond lengths; check chemical accuracy. Use 8q $\mathrm{H}_2\mathrm{O}$
(contextual) as a negative control: the simulator should fail or scale badly.

\subsection*{Q3. How does the contextuality verdict correlate with alternative advantage-witnesses (magic, entanglement, Wigner negativity) on the same six Table~I Hamiltonians?}
\textbf{Basis.} The paper implicitly claims contextuality is \emph{the}
non-classicality witness for VQE. Magic (stabilizer rank), entanglement
entropy, and Wigner negativity give correlated but not identical verdicts;
disagreements are scientifically informative because contextuality is an
operator/measurement-level notion, magic is state-level.\\
\textbf{Next steps.} For each Table~I ground state: compute Renyi-2 entropy
on a symmetric bipartition; compute stabilizer-rank lower bound via
Bravyi-Gosset for $\le 8$ qubits; compute discrete Wigner negativity where
applicable. Build a $6$-row cross-tabulation.

\subsection*{Q4. What is the noise threshold at which a nominally-contextual VQE run becomes indistinguishable from a noncontextual one under realistic hardware noise + finite shots?}
\textbf{Basis.} The classifier operates on ideal Pauli operators. On NISQ
hardware, expectation values are measured with depolarizing/dephasing noise
and shot noise; boundary cases like $\mathrm{HeH}^+$ ($CD_0{=}0.38$) could
lose their witness triples under realistic error rates.\\
\textbf{Next steps.} Simulate noisy shot-based estimation (Qiskit Aer;
1q err $10^{-3}$, 2q err $10^{-2}$, readout $10^{-2}$) on
$\mathrm{HeH}^+$ 4q. For shot budgets $\{10^3, 10^4, 10^5, 10^6\}$, plot the
probability of recovering a valid Theorem~3 witness. Use LiH and
$\mathrm{H}_2\mathrm{O}$ as deeper-in-contextual-regime controls.

\subsection*{Q5. Does Theorem~3 extend meaningfully to QAOA/AQC, or does it produce trivially non-contextual verdicts for canonical QAOA (Max-Cut) because the cost Hamiltonian is diagonal in $Z$?}
\textbf{Basis.} Max-Cut / MaxSAT QAOA cost Hamiltonians are entirely
Z-diagonal, hence mutually commuting, hence $T = \emptyset$ and trivially
non-contextual --- which contradicts the intuition that some QAOA instances
have quantum-speedup potential. Either QAOA truly has no contextuality-based
advantage, or the framework needs an extension (e.g., include the mixer
Hamiltonian in the operator set).\\
\textbf{Next steps.} (1)~Confirm the trivial-non-contextual verdict for
Max-Cut QAOA on 3--6 test graphs. (2)~Extend the operator set to include
mixer Pauli terms; re-run and check whether the combined set becomes
contextual. (3)~For AQC, apply to $H(s)$ at $s\in\{0, 0.5, 1\}$. (4)~Publish
even a negative finding as a QAOA-advantage-literature clarification.
